\subsection{Reconstruction of survival of one prescribed finite quotient}

\begin{result}{E.20}{Proved}{SECTIONS\_OVER\_Q A.72 and GT\_TOWER\_CHECK C.2.17 confirmed}
Let \(N\subset\Pi\) be open and normal in \(E_U\), and put
\(Q=\Pi/N\). There is an integer \(N_0>0\) and an ordinary ghost
\[
 A=3\prod_{j\ge0}q_j^{N_0(j+4)!},\qquad
 B=-2\prod_{j\ge0}r_j^{N_0(j+4)!}
\]
whose section in
\[
 E_U/(N\cap\overline{[\Pi,\Pi]})
 \simeq E_Q\mathbin{\times}_{E_{Q^{\rm ab}}}E_{\rm ab}
\]
has \(Q\)-component the specified point-three section. At every finite
place this entire section is geometrically conjugate to the image of the
ordinary local point of the pair. The pair and the section depend on \(N\).
Consequently no one fixed finite geometric quotient excludes every A.72
output. This last assertion confirms R7, with its varying-pair quantifier.
\end{result}
\begin{proof}
We reconstruct the arithmetic and compatibility argument, retaining the
full abelian factor. Put \(h=|Q|\), \(e=\exp Q^{\rm ab}\), and
\(\Lambda=\Pi^{\rm ab}=\widehat{\mathbb Z}(1)^2\).
The subgroup \(Ns_3(G_{\mathbb Q})\) defines a finite etale cover
\(Y\to U\) of degree \(h\) with a rational point \(y_0\) above three.
Its geometric monodromy is the regular \(Q\)-action.

Choose a finite Galois field \(F/\mathbb Q\) over which all its boundary
points, deck transformations, and leading coefficients at zero and one
are defined. Include \(\mu_h,3^{1/h},(-2)^{1/h}\). At each boundary
point of index \(a\mid h\), write the map as
\(t=d z^a(1+zh(z))\), or use \(1-t\), and include \(d^{1/a}\)
in \(F\). An integer \(\Delta\), divisible by \(6h\), makes these
models and parameters integral and smooth away from its divisors.
At a prime split in \(F\) outside \(\Delta\), taking the \(a\)-th
root of \(1+zh(z)\) is an invertible analytic coordinate change on the
maximal ideal. Every positive-valuation \(a\)-th power value on that
branch therefore lifts, with all \(a\) local roots. When this holds at
every boundary point, the whole fibre of \(Y\) is rational.

Here is the field-intersection calculation needed for the prime choices.
Let \(M=[F:\mathbb Q]\), and put
\(K_d=F(\mu_{4M},d^{1/(4M)})\) for \(d=3,-2\).
By Kronecker--Weber, each \(K_d^{\rm ab}\) is contained in
\(\mathbb Q(\mu_{f_d})\) for an integer \(f_d\)
\citep[I, Theorem 4.16, pp.~49--50]{RootMilneCFT}.
Choose \(C\) divisible by \(24,4M,h\), and \(N_0\) divisible by
\(C,\Delta,f_3,f_{-2}\). For \(n\) divisible by these parameters and
\(d'=dw^C\), with positive rational \(w\) a unit at two and three,
let \(L=\mathbb Q(\mu_n,(d')^{1/n})\).
The radical calculation in E.102, applied to \(w^{C/24}\), gives
\[
 H=\operatorname{Gal}(L/\mathbb Q(\mu_n))\simeq\mathbb Z/(n/2).
\]
The square root is already cyclotomic. Any further odd-prime root or
fourth root would put a nonnormal pure radical field inside an abelian
extension of \(\mathbb Q\), which the valuation-one argument of E.102
excludes. This explains both the order and the exceptional factor two.
The subgroup
\(H_F=\operatorname{Gal}(FL/F\mathbb Q(\mu_n))\subset H\)
has index dividing \(M\), so \(MH\subset H_F\).
Complex conjugation in \(\operatorname{Gal}(FL/\mathbb Q)\) acts on
\(H_F\) by inversion. Hence its commutator subgroup contains \(2MH\).
The fixed field of this last subgroup is
\[
 F\mathbb Q(\mu_n,(d')^{1/(4M)})=K_d\mathbb Q(\mu_n).
\]
For a finite Galois field \(K\) and an abelian field \(T\), the Galois
fibre product for \(KT\) has commutator subgroup
\([\operatorname{Gal}(K/\mathbb Q),\operatorname{Gal}(K/\mathbb Q)]
\times\{1\}\): every element of the first factor lifts, and pairwise
commutators are calculated coordinatewise. Thus
\((KT)^{\rm ab}=K^{\rm ab}T\). Using the chosen conductors gives
\[
 (FL)^{\rm ab}=\mathbb Q(\mu_n),\qquad
 FL\cap\mathbb Q(\mu_P)=\mathbb Q(\mu_n)\quad(n\mid P).
\]
A unit residue class modulo \(P\) that is one modulo \(n\) is therefore
compatible with complete splitting in \(FL\). Chebotarev
\citep[Theorem 8.31, p.~147]{CohMilneANT} supplies arbitrarily large
primes with both properties. The specified pair of restrictions is
central in the relevant Galois fibre product, so its conjugacy class
has the same restrictions; no choice of incompatible representatives
is hidden in this use of Chebotarev.

For \(a=3u^C,b=-2v^C\), define the diagonal cover
\(V_{a,b,n}:aX^n+bY^n=1\), with \(XY\ne0\), and let
\(Z_Q=Y/[Q,Q]\). Since \(e\mid C,n\), the coordinates
\[
 (a/3)^{1/e}X^{n/e},\qquad (b/(-2))^{1/e}Y^{n/e}
\]
give a rational morphism \(V_{a,b,n}\to Z_Q\), through the pointed
\(e\)-Kummer cover. The largest common geometric quotient of \(Y\)
and \(V_{a,b,n}\) is exactly \(Q^{\rm ab}\): a common quotient is
abelian and the divisibility of \(n\) makes all of \(Q^{\rm ab}\)
available. Thus
\[
 W_{a,b,n}=Y\times_{Z_Q}V_{a,b,n}
\]
is geometrically connected and has group
\(Q\times_{Q^{\rm ab}}\Lambda/n\Lambda\), of degree at most
\(D=hn^2\). Its compactification has \(2g\le D+2\) and at most
\(3D\) boundary points, by Riemann--Hurwitz.
At \(p\nmid\Delta nab\), prime-to-\(p\) specialization preserves
this connected cover and its boundary inertia
\citep[Expose XIII, Corollary 2.9, section 2.10, Lemma 2.11,
pp.~288--290]{RootSGA1}. The local tame normalizations are smooth:
after an unramified extension their equations reduce to the power
parameters above, whose exponents are prime to \(p\).
The trace formula and Weil bound
\citep[(1.5.1) and Theorem (1.6), pp.~275--276]{RootDeligne1974}
give at least \(p+1-(D+2)\sqrt p-3D\) interior points.
This is positive for \(p>B_Q(n)=(4hn^2+10)^2\), since at the endpoint
the expression equals \(12D^2+59D+81>0\) and is increasing thereafter.
Smooth Hensel lifting supplies a local point with unit \(X,Y\).

Set \(n_j=N_0(j+4)!\), \((a_0,b_0)=(3,-2)\), and
\(W_j=W_{a_j,b_j,n_j}\). We maintain unit-coordinate local points
at every finite prime. At \(p\le B_Q(n_j)\) and \(p\mid\Delta\)
we maintain a point over \(t=a_j\), with \(X=1\), and the bounds
\[
 Y\equiv1\pmod{p^{v_p(n_{j+1})}},\qquad
 v_p(a_j+b_j-1)\ge V_j(p)+2v_p(n_{j+1})+j+3,
\]
where \(V_j(p)=\max(v_p(a_j),v_p(b_j))\).
The rational point \((y_0,1,1)\) starts the induction.
At stage \(j\), put \(n=n_j,m=n_{j+1}\), and let \(S\) contain
all primes at most \(B_Q(m)\), the divisors of \(\Delta\), and all
coefficient-support primes. Choose a unit-coordinate point
\((y_p,\xi_p,\eta_p)\) of \(W_j\) at each \(p\in S\).
For \(p\mid m\), use the maintained point; indeed such a prime is
a divisor of \(N_0\) or is at most \(j+5\), and is already below
\(B_Q(n_j)\). Thus \(\xi_p,\eta_p\equiv1\pmod{p^{v_p(m)}}\).

Choose
\(L_p\ge V_j(p)+2v_p(n_{j+2})+j+4\), and increase these finitely
many precisions as needed below. If units \(Q_p,R_p\) have these
residues \(\xi_p,\eta_p\), the etale inverse of \(Y\to U\)
gives a point \(y'_p\) near \(y_p\) over \(a_jQ_p^n\).
Newton's criterion at one for
\[
 Z^m-\frac{1-a_jQ_p^n}{b_jR_p^n}
\]
applies, because its error has valuation at least \(L_p-v_p(b_j)>2v_p(m)\).
It gives \(Z_p\) with
\(v_p(Z_p-1)\ge L_p-v_p(b_j)-v_p(m)\), hence
\(Z_p\equiv1\pmod{p^{v_p(n_{j+2})}}\).
The point \((y'_p,Q_p,R_pZ_p^{m/n})\) belongs to the selected component
of the old fibre product after increasing \(L_p\): it tends to the
original point, and the component is open and closed. This is a
characteristic-zero etale assertion, also when \(p\mid e\).

The preceding field-intersection lemma now chooses a new prime \(q_j\)
larger than all the retained primes, with residues \(\xi_p\) modulo
\(p^{L_p}\), and split in
\(F\mathbb Q(\mu_n,b_j^{1/n})\).
At \(q_j\), simple-root Hensel lifting gives a unit \(z\) with
\(b_jz^n=1-a_jq_j^n\).
It then chooses a new prime \(r_j>q_j\), with residues \(\eta_p\)
at the old primes, residue \(z\) modulo \(q_j\), and split in
\(F\mathbb Q(\mu_m,(a_jq_j^n)^{1/m})\).
All CRT classes are one modulo the required exponent; the new radicand
is still three times a rational \(C\)-th power. Both applications
therefore satisfy the lemma's hypotheses.
Set \(a_{j+1}=a_jq_j^n,b_{j+1}=b_jr_j^n\).
The rational transition morphism is
\[
 W_{j+1}\longrightarrow W_j,\qquad
 (y,X,Y)\longmapsto(y,q_jX^{m/n},r_jY^{m/n}).
\]
It preserves the two displayed common abelian coordinates, hence the
specified connected component. The Newton calculation gives the next
maintained point \((y'_p,1,Z_p)\) at each old \(p\), together with
the required sum-error bound. At \(q_j\), the residue chosen for
\(r_j\) makes \((1-a_jq_j^n)/(b_jr_j^n)\) an \(m\)-th power.
The value \(t=a_jq_j^n\) is an appropriate positive-valuation power
at every zero boundary: use \(h\mid C,n\) and the radical in \(F\).
The whole \(Y\)-fibre therefore splits, so the desired abelian sheet
has a lift. At \(r_j\) interchange the two coordinates and use the
one boundary and its splitting field. The remaining primes have unit
coefficients, do not divide \(\Delta m\), and exceed \(B_Q(m)\);
the Weil bound gives their unit-coordinate points. These exhaust the
local cases, including two and three.

The products now converge in the power completion, since \(n_j/C\)
is eventually divisible by every fixed integer. Literal relative Kummer
cocycles \(\sum_jn_j\kappa(q_j),\sum_jn_j\kappa(r_j)\) converge
uniformly and have zero image in \(Q^{\rm ab}\). They combine with
the fixed literal \(Q\)-section \(s_3\) to give a continuous section
in the stated fibre product. The group isomorphism in the statement
comes from the diagonal map: its kernel is the displayed intersection,
and adjusting two representatives by their common abelian difference
proves surjectivity. The relevant product of normal subgroups is closed
by compactness, so this is an isomorphism of profinite groups over \(G\).

At a fixed prime \(p\), valuations of \(a_j,b_j\) stabilize. Their
unit factors converge ordinarily, since \(n_j\) eventually kills every
finite quotient of \(\mathbb Z_p^\times\). Their limits \(x_p,y_p\)
are nonzero, and the maintained sum-error gives \(x_p+y_p=1\).
Their finite-power classes are the localizations of \(A,B\).
For a point on \(W_j\) with unit coordinates, its image
\(a_jX_j^{n_j}\) tends to \(x_p\), uniformly in the possible units.
For fixed \(i\) every such image with \(j\ge i\) lifts to \(W_i\).
The image of \(W_i(\mathbb Q_p)\) in \(U(\mathbb Q_p)\) is closed:
use compactness of the projective compactification and
\(x_p\ne0,1,\infty\). Hence \(x_p\) lifts to every \(W_i\).
The corresponding finite geometric conjugator sets are nonempty and
nested. Their inverse images in the compact geometric group of the
fibre product have a common point, which is one conjugator for the
whole local section. Positive roots of the added primes make its real
restriction the image of three. Every chosen support prime retains its
positive ordinary valuation, so both supports are infinite. E.101
therefore excludes a rational pair. Projection proves the R7 conclusion.
\end{proof}

The independent certificate \path{code/root_finite_quotient_checks.py}
checks the cyclic index containment, the point-bound polynomial, the
common-coordinate exponent identity, and actual Newton lifts with
nonunit derivatives. These finite checks support the indicated algebra;
they do not replace the field and curve arguments above.
